\setcounter{section}{1}
\setcounter{theorem}{9}
\section{Preliminaries}\label{sec: prelims}
For any positive integer $n$, we denote the set $\{1, \dots, n\}$ by $[n]$. For $d \leq n$ we use the notation $\binom{n}{\leq d} = \binom{n}{0} + \cdots + \binom{n}{d}$. Note that $\binom{n}{\leq d} < (n+1)^d$.

In this section we review the necessary preliminaries.  We first review some basics of Fourier analysis on the Boolean cube.
Consider the vector space of functions from $\pmone^n$ to $\R$, equipped with an inner product defined by
\[
\langle f, g \rangle := \E_{x \in \pmone^n}[f(x)g(x)] = \frac{1}{2^n}\sum_{x \in \pmone^n}f(x)g(x).
\]
for every $f,g: \pmone^n \rightarrow \R$. For any set $S \subseteq [n]$, define the associated \emph{parity} function $\chi_S$ by $\chi_S(x) = \prod_{i \in S}x_i$. The set of parity functions $\bra{\chi_S : S \subseteq [n]}$ form an orthonormal basis for this vector space.  Thus, every function $f : \pmone^n \to \R$ has a unique multilinear expression as
\[
f = \sum_{S \subseteq [n]}\wh{f}(S)\chi_S.
\]
We also use $\PARITY_n$ to denote $\chi_{[n]}:\pmone^n\to\pmone$.
The coefficients $\bra{\wh{f}(S) : S \subseteq [n]}$ are called the \emph{Fourier coefficients} of $f$.
\begin{fact}[Parseval's Identity]\label{fact: parseval}
For any function $f : \pmone^n \to \R$,
\[
    \sum_{S \subseteq [n]}\wh{f}(S)^2 = \frac{\sum_{x \in \pmone^n}f(x)^2}{2^n}.
\]
\end{fact}

\begin{definition}[Spectral Norm]\label{defn: wt}
For any function $f : \pmone^n \to \R$, define its \emph{spectral norm}, which we denote $\snorm{f}$, to be the sum of absolute values of the Fourier coefficients of $f$.  That is,
\[
\snorm{f} : = \sum_{S \subseteq [n]}\abs{\wh{f}(S)}.
\]
\end{definition}

\begin{definition}[Hadamard Codeword]\label{defn: hcwd}
If a $\ell$-bit string $(x_1, \dots, x_{\ell}) \in \pmone^\ell$ (alternatively, view the indices of $x$ as subsets of $[\log\ell]$) is of the form $x_S = \prod_{i \in S}z_i$ for all $S \subseteq [\log \ell]$ for some $z \in \pmone^{\log \ell}$, then define such an $x = x_1 \dots x_\ell$ to be the \emph{$\ell$-bit Hadamard codeword} $h(z)$ of the $(\log \ell)$-bit string $z$.
\end{definition}

\subsection{Addressing functions}
\begin{definition}[$(m,k)$-addressing function]\label{defn: add}
We define a (partial) function $f:\pmone^{m+k} \rightarrow \bra{-1, 1, \star}$ to be an \emph{$(m,k)$-addressing function} if there exists $g : \pmone^m \to \bra{[k] \cup \star}$ such that 
\begin{itemize}
    \item $f(x_1, \dots, x_m, y_1, \dots, y_{k}) = y_{g(x_1, \dots, x_m)}$ if $g(x_1, \dots, x_m) \in [k]$, and $f(x_1, \dots, x_m, y_1, \dots, y_{k}) = \star$ otherwise.
    \item For all $j \in [k]$, there exists $(x_1, \dots, x_m) \in \pmone^m$ such that $g(x_1, \dots, x_m) = j$.
\end{itemize}

We call the variables $\{x_1, \dots, x_m\}$ the \emph{address variables} and the variables $\{y_{1}, \dots, y_{k}\}$ the \emph{target variables}.  The function $g$ is called the $\what$ function of $f$.
\end{definition}

\begin{definition}[Indexing Function]
The Indexing function, which we denote by $\IND_k$, is a $(k, 2^k)$-addressing function defined by $\IND(x_1, \dots, x_k, y_1, \dots, y_{2^k}) = y_{\textnormal{bin}(x)}$, where $\textnormal{bin}(x)$ denotes the integer represented by the binary string $x_1, \dots, x_k$.
\end{definition}

\begin{definition}[Composition with addressing functions]\label{defn: addcomp}
For any function $f : \pmone^n \to \pmone$ and an $(m, k)$-addressing function $\ADD$, define the (partial) function $f^{\ADD} : \pmone^{n (m + k)} \to \bra{-1, 1, \star}$ by
\[
f^{\ADD}(x_1, y_1, \dots, x_n, y_n) = 
\begin{cases}
f(\ADD(x_1, y_1), \dots, \ADD(x_n, y_n)) & \text{if~} \ADD(x_i, y_i) \in \pmone \text{~for all~} i \in [n]\\
\star & \text{otherwise}.
\end{cases}
\]
where $x_i \in \pmone^{m}$ and $y_i \in \pmone^k$ for all $i \in [n]$.
\end{definition}

\begin{definition}[Hadamard Addressing Function]\label{defn: hadd}
We define the \emph{Hadamard addressing function}, which we denote $\HADD_\ell : \pmone^{2\ell} \to \bra{-1, 1, \star}$, as follows.
Fix an arbitrary order on the $\ell$-many Hadamard codewords of $(\log \ell)$-bit strings, say $w_1, \dots, w_\ell$.  Define the selector function of $\HADD_\ell$ by
\begin{align*}
g(x) & =\begin{cases}
i & \text{if~} x = w_i \text{~for some~} i \in [\ell]\\
\star & \text{otherwise}.
\end{cases}
\end{align*}
Note that $\HADD_\ell$ is an $(\ell, \ell)$-addressing function.
\end{definition}

\subsection{Polynomial approximation}

\begin{definition}[Approximate Degree]\label{defn: adeg}
The \emph{$\eps$-approximate degree} of a function $f : \pmone^n \to \bra{-1, 1, \star}$, denoted by $\adeg_\eps(f)$ is defined to be the minimum degree of a real polynomial $p : \pmone^n \to \R$ that satisfies $\abs{p(x) - f(x)} \leq \eps$ for all $x \in \pmone^n$ for which $f(x) \in \pmone$.\footnote{When dealing with partial functions, another notion of approximation is sometimes considered, where the approximating polynomial $p$ is required to have bounded values even on the non-promise inputs of $f$.  For the purpose of this paper, we do not require this constraint.} That is,
\[
\adeg_\eps(f):=\min\left\{d:\begin{array}{l}\deg(p)\leq d,\quad |p(x)-f(x)|\leq\epsilon\\\text{for all }x\in\pmone^n\text{ for which }f(x)\in\pmone\end{array}\right\}.
\]
\end{definition}
Henceforth, we will use the notation $\adeg(f)$ to denote $\adeg_{1/3}(f)$.

\begin{definition}[Approximate Spectral Norm]\label{defn: awt}
The \emph{approximate spectral norm} of a function  $f : \pmone^n \to \{-1,1,\star\}$, denoted by $\asnorm{f}{\epsilon}$ is defined to be the minimum spectral norm of a real polynomial $p : \pmone^n \to \R$ that satisfies $\abs{p(x) - f(x)} \leq \eps$ for all $x \in \pmone^n$ for which $f(x) \in \pmone$.
\[
\asnorm{f}{\epsilon}:=\min\left\{\snorm{p}:\begin{array}{l}|p(x)-f(x)|\leq\epsilon\\\text{for all }x\in\pmone^n\text{ for which }f(x)\in\pmone\end{array}\right\}.
\]
\end{definition}
\begin{lemma}[\cite{BNRW07}]\label{lem: asnormequiv} 
Let $f:\pmone^n \to \pmone$ be a total function. Then for all constants $0 < \delta, \eps < 1$ we have $\adeg_\eps(f) = \Theta(\adeg_\delta(f))$. The constant hidden in the $\Theta$ notation depends on $\delta$ and $\eps$.
\end{lemma}

The following is a standard upper bound on the approximate spectral norm of a Boolean function in terms of its approximate degree.
\begin{claim}\label{claim: adegnaivelb}
For all total functions $f : \pmone^n \to \pmone$, we have 
\[
\log \asnorm{f}{1/3} = O(\adeg(f) \log n).
\]
\end{claim}
\begin{proof}
Let $d$ denote the approximate degree of $f$.  Take any 1/3-approximating polynomial of degree $d$, say $p$, to $f$.  Then,
\begin{align*}
    \sum_{S \subseteq [n]} \abs{\wh{p}(S)} & \leq \sqrt{\binom{n}{\leq d}} \cdot \sqrt{\sum_{S: |S| \leq d}\wh{p}(S)^2} \leq 4/3 \cdot (n+1)^{d/2} = 2^{O(d \log n)},
\end{align*}

where the first inequality follows by the Cauchy-Schwarz inequality, the second inequality follows by Parseval's identity (Fact~\ref{fact: parseval}) and the fact that the absolute value of $p$ is at most $4/3$ for any input $x \in \pmone^n$.
\end{proof}

It is easy to exhibit functions $f : \pmone^n \to \pmone$ such that $\log \asnorm{f}{1/3} = \Omega(\adeg(f))$.  Bent functions satisfy this bound, for example.

Building upon ideas in~\cite{KP97}, the approximate spectral norm of $f \circ \IND_1$ was shown to be bounded below by $2^{\Omega(\adeg(f))}$ in~\cite{CM17}.
\begin{theorem}[\cite{CM17}]
Let $f:\pmone^n \rightarrow \pmone$ be any function. Then $\asnorm{f\circ\IND_1}{1/3} \geq 2^{c \cdot \adeg_{2/3}(f)}$ for any constant $c < 1 - 3/\adeg_{2/3}(f)$.
\end{theorem}

\subsection{Communication complexity}
The classical model of communication complexity was introduced by Yao in~\cite{Yao79}. In this model two parties, say Alice and Bob, wish to compute a function whose output depends on both their inputs. Alice is given an input $x \in \mathcal{X}$, Bob is given $y \in \mathcal{Y}$, and they want to jointly compute the value of a given function $F(x,y)$ by communicating with each other.  Alice and Bob individually have unbounded computational power and the number of bits communicated is the resource we wish to minimize. Alice and Bob communicate using a \emph{protocol} that is agreed upon in advance. In the randomized model, Alice and Bob have access to unlimited public random bits and the goal is to compute the correct value of $F(x,y)$ with probability at least $2/3$ for all inputs $(x,y) \in \mathcal{X} \times \mathcal{Y}$. The \emph{bounded-error randomized communication complexity} of a function $F$, denoted $R^{cc}(F)$, is the number of bits that must be communicated in the worst case by any randomized protocol to compute the correct value of the function $F(x,y)$, with probability at least $2/3$, for every $(x,y) \in \mathcal{X} \times \mathcal{Y}$.

The quantum model of communication complexity was introduced by Yao in~\cite{Yao93}. 
We refer the reader to the survey~\cite{Wolf02} for details. The \emph{bounded-error quantum communication complexity} of a function $F$, denoted $Q^{cc}(F)$ is the number of bits that must be communicated by any quantum communication protocol in the worst case to compute the correct value of the function $F(x,y)$, with probability at least $2/3$, for every $(x,y)$ in domain of $F$. Buhrman, Cleve and Wigderson~\cite{BCW98} observed a quantum simulation theorem, which gives an upper bound on the bounded-error quantum communication complexity of a composed function of the form $f \circ \wedge$ or $f \circ \oplus$ in terms of the bounded-error quantum query complexity of $f$ (see Theorem~\ref{thm: qtm_sim_log_loss}).

Lee and Shraibman~\cite{LS09} showed that the bounded-error quantum communication complexity of $f \circ \oplus$ is bounded below by the logarithm of the approximate spectral norm of $f$. Also see~\cite{CM17} for an alternate proof.
\begin{theorem}[\cite{LS09}]\label{thm: LS}
For any Boolean function $f : \pmone^n \to \pmone$,
\[
Q^{cc}(f \circ \oplus) = \Omega(\log \asnorm{f}{1/3}).
\]
\end{theorem}



